\documentclass[10pt]{article}
\usepackage[margin=0.7in]{geometry}
\usepackage{enumitem}
\usepackage{sectsty}

\sectionfont{\large\bfseries}
\subsectionfont{\normalsize\bfseries}

\title{\vspace{-2em}\Large \textbf{Project Defense Cheat Sheet: Face Recognition Attendance System}}
\author{}
\date{\vspace{-2em}}

\begin{document}
\maketitle

\section*{1. Elevator Pitch}
An automated, local attendance system that replaces manual roll-calls with real-time facial recognition. It provides a web-based dashboard (React/Next.js) for educators to register students, train the recognition model, and mark attendance automatically via webcam using a Python (FastAPI) backend.

\section*{2. Technology Stack \& Justifications (Why did you choose these?)}
\begin{itemize}[leftmargin=*, parsep=0pt, itemsep=3pt]
    \item \textbf{Backend:} Python / FastAPI. \textit{Why?} High performance, excellent integration with data/ML libraries (OpenCV), and automatic API documentation.
    \item \textbf{Computer Vision:} OpenCV (Haar Cascades \& LBPH). \textit{Why?} Efficient enough to run on standard CPUs without requiring expensive GPUs. 
    \item \textbf{Frontend:} Next.js \& Tailwind CSS. \textit{Why?} Rapid UI development, responsive modern interface, and component-based architecture.
    \item \textbf{Database:} SQLite (via SQLAlchemy). \textit{Why?} Zero-configuration, file-based database ideal for an academic, standalone project.
\end{itemize}

\section*{3. How It Works (The Workflow)}
\begin{enumerate}[leftmargin=*, parsep=0pt, itemsep=3pt]
    \item \textbf{Registration:} UI sends student details. Backend opens camera, detects face using Haar Cascade, captures 30 grayscale face crops, and saves them locally.
    \item \textbf{Training:} The LBPH recognizer processes the saved dataset, extracting features to create a single trained model file (\texttt{trainer.yml}).
    \item \textbf{Recognition \& Attendance:} During live feed, faces are continuously detected. The LBPH model predicts the ID based on confidence scores. Recognized students are logged in SQLite (strictly once per day).
\end{enumerate}

\section*{4. Key Algorithms (Crucial for VIVA/Review)}
\begin{itemize}[leftmargin=*, parsep=0pt, itemsep=3pt]
    \item \textbf{Face Detection (Haar Cascades):} An object detection method based on Haar-like features. It slides a window over the image, quickly discarding non-face regions. \textit{Optimization used:} We scale frames down by 50\% during detection to maintain high FPS.
    \item \textbf{Face Recognition (LBPH):} Local Binary Patterns Histograms. Divides the face into small grids, compares each pixel to its neighbors to create a binary number, forming a histogram. It compares these histograms to known records. \textit{Benefit:} Works well even if lighting varies slightly.
\end{itemize}

\section*{5. Known Limitations \& Future Scope (Anticipate their critiques)}
\begin{itemize}[leftmargin=*, parsep=0pt, itemsep=3pt]
    \item \textbf{Spoofing Vulnerability:} Can be tricked by a photograph. \textit{Future Fix:} Integrate liveness detection (e.g., blink detection or depth sensing).
    \item \textbf{Security:} No login system for the administrator dashboard. \textit{Future Fix:} Add JWT-based authentication.
    \item \textbf{Scalability:} SQLite and LBPH might slow down with thousands of students. \textit{Future Fix:} Migrate to PostgreSQL and a Deep Learning model (e.g., FaceNet, dlib) for large-scale deployments.
    \item \textbf{Lighting Dependency:} Heavy shadows can lower accuracy. \textit{Future Fix:} Add pre-processing steps like histogram equalization to normalize lighting.
\end{itemize}

\end{document}
